% §5.4 One-Sided vs Two-Sided Approximate Sets
%
% Cross-reference resolution status:
%
% DEFINED here:
%   \label{sec:one-sided-two-sided}             -- this section
%   \label{def:one-sided}                       -- one-sided approximate set
%
% RESOLVED (already in source tree from Plans 1-4 and Plan 5 Tasks 2-4):
%   \Cref{sec:from-atom-to-set}                 -- chapter 5 §5.1
%   \Cref{sec:bloom-from-scratch}               -- chapter 5 §5.2
%   \Cref{def:bloom-filter}                     -- chapter 5 §5.2
%   \Cref{sec:predicate-algebra}                -- chapter 5 §5.3
%   \Cref{sec:three-faces}                      -- chapter 1 §1.1
%   \Cref{sec:asymmetric-channel}               -- chapter 2 §2.2
%   \Cref{def:z-channel}                        -- chapter 2 §2.2
%   \Cref{sec:bridge-to-sets}                   -- chapter 3 §3.5
%   \Cref{ch:set-composition}                   -- chapter 6 stub label
%   \Cref{ch:perfect-hash-filters}              -- chapter 15 stub label
%   \Cref{ch:optimal-bernoulli-sets}            -- chapter 14 stub label
%
% FORWARD REFS within chapter 5 (resolve in Plan 5 Task 6):
%   \Cref{sec:ch5-what-comes-next}              -- §5.5 (Task 6)
%
% Section-level forward references into chapters 6, 14, 15 do NOT exist yet.
% The CHAPTER labels above resolve via the Plan 0 stubs; sub-section labels
% inside those chapters are unresolved until Plans 6+ land. Where this matters
% (the perfect-filter line in chapter 6 and the BPHF in chapter 15), we use
% chapter-level \Cref or prose.

\section{One-Sided vs Two-Sided Approximate Sets}
\label{sec:one-sided-two-sided}

The Bloom filter of \Cref{sec:bloom-from-scratch} has a striking
property: one of its two error rates is exactly zero, not by parameter
choice but by construction. The bit array can only accumulate ones,
so an inserted item is never reported missing; the other rate is
positive and tunable. This is one shape an approximate set can take.
A different construction might force the false-positive rate to zero
on a designated set instead, leaving the false-negative rate positive;
a third might allow both rates to be positive. The present section
gives the three shapes their names and connects the distinction to
the algebraic structure that chapter~6 will develop.

\subsection*{Definition}

\begin{definition}[One-sided and two-sided approximate sets]
\label{def:one-sided}
Let $\ASet{A} : \Universe \to \Bernoulli{\BoolType}$ be an approximate
set of $A \subseteq \Universe$ with per-block rates
$(\fprate, \fnrate)$ as in \Cref{sec:predicate-algebra}. We say
$\ASet{A}$ is
\begin{itemize}
  \item \emph{one-sided in FNR} (or \emph{FNR-zero}) if $\fnrate = 0$
    and $\fprate > 0$;
  \item \emph{one-sided in FPR} (or \emph{FPR-zero}) if $\fprate = 0$
    and $\fnrate > 0$;
  \item \emph{two-sided} if both $\fprate > 0$ and $\fnrate > 0$.
\end{itemize}
The Bloom filter of \Cref{def:bloom-filter} is one-sided in FNR. The
perfect hash filter of \Cref{ch:perfect-hash-filters} is one-sided in
FPR on its design domain.
\end{definition}

The names track the rate that is zero, not the one that is positive.
The Z-channel of \Cref{def:z-channel} is the per-query analog; the
present section's one-sided property is its per-set lift. A two-sided
approximate set sits in the interior of the rate square $[0,1]^2$; a
one-sided set sits on one of the two coordinate axes. The corner
$\fprate = \fnrate = 0$ is the exact set, not an approximate set at
all, and is excluded from \Cref{def:one-sided}.

\subsection*{Why FNR-zero is the natural default for insertion-only structures}

The Bloom filter's FNR-zero property was structural, not designed. The
insert operation sets bits to $1$, and no operation ever resets them;
a subsequent query for an inserted item reads bits that were set, so
the query returns $\top$ regardless of what else has happened since.
The argument depended only on the monotonicity of the bit array.
Insertions go up, never down. Any data structure built on the same
monotone-insert primitive shares the property: counting sketches keyed
by hash, summary structures over append-only logs, and any
representation of an insertion-only bit-set query (read what was
written) are one-sided in FNR for the same structural reason.

The asymmetry runs the other way. Forcing $\fprate = 0$ on a designated
set, while keeping $\fnrate > 0$ elsewhere, requires the construction
to know which queries are guaranteed inside and which are not. A
deterministic check against the underlying set would do this, but at
the cost of storing the set itself, which destroys the space saving
an approximate set was supposed to provide. The constructive path to
an FPR-zero structure must be more clever. Chapter~15's perfect hash
filter (\Cref{ch:perfect-hash-filters}) computes a perfect hash for a
fixed $A$ at construction time and stores only a fingerprint per
element; a query for $x \in A$ lands on its fingerprint and is
reported correctly with probability $1$, while a query for
$x \notin A$ may still collide. The set has to be known up front, and
updates are expensive. Those costs buy the opposite asymmetry:
certainty in the positive direction rather than the negative.

\subsection*{The trade-off: safety versus space}

One-sided approximate sets are valuable because one direction of the
answer is exact. If a Bloom filter reports $\bot$, the queried element
is definitely not in the underlying set; downstream logic can act on
the certain answer without further checks. An FPR-zero filter has the
dual safety: a $\top$ report is certain, only $\bot$ needs follow-up.
Many systems are built around exactly this asymmetric trust. Web
caches use a Bloom filter to skip lookups for keys known to be absent;
the certain $\bot$ avoids a slow disk read. Spell-checkers flag
possibly misspelled words; the certain $\top$ on a known word avoids
a distracting underline. In every such setting, the one-sided
structure substitutes for an exact set precisely because the certain
direction is the one the application needs.

Two-sided sets give up the certain direction and recover space.
Counting Bloom filters store small integer counters instead of bits
to support deletions, paying for that with a nonzero $\fnrate$.
Cardinality estimators such as HyperLogLog and generic sublinear
sketches are typically two-sided on derived predicates. The trade is
honest: dropping one side of the asymmetry frees the construction to
use space more aggressively.

A worked comparison makes the size of the trade visible. A Bloom
filter with target $\fprate = 0.01$ and $\fnrate = 0$ requires
$m \approx 10\, n \log_2 e \approx 14.4\, n$ bits, by the optimal-$k$
formula of \Cref{sec:bloom-from-scratch} solved for $m/n$ at the
given $\fprate$. Suppose a hypothetical two-sided set achieves
$\fprate = \fnrate = 0.05$ at $m \approx 7\, n$ bits. On a balanced
workload of $1000$ in-set and $1000$ out-of-set queries, the Bloom
filter expects $0.01 \cdot 1000 = 10$ false positives and $0$ false
negatives. The two-sided set expects $0.05 \cdot 1000 + 0.05 \cdot
1000 = 100$ misclassifications, ten times more, but in roughly half
the space. The hypothetical framing is honest: chapter~14's Bernoulli
Hash Function is the real space-optimal one-sided construction, and
actual two-sided constructions that beat the Bloom filter are more
constrained than the comparison suggests. The point is the qualitative
shape. One-sidedness is paid for in bits, and a less safety-conscious
construction can in principle spend those bits elsewhere.

\subsection*{The perfect filter line}

The two-dimensional rate plane $(\fprate, \fnrate) \in [0, 1]^2$
parameterizes every approximate set. The one-sided sets occupy the
two coordinate axes: the FNR-zero family on
$\{\fnrate = 0,\, 0 < \fprate \le 1\}$ and the FPR-zero family on
$\{\fprate = 0,\, 0 < \fnrate \le 1\}$. We call the union of these
two segments the \emph{perfect filter line}. The name is the
convention of the underlying paper family; the word ``line'' is
appropriate because each segment is a one-dimensional locus inside
the two-dimensional plane. A one-sided approximate set has one degree
of rate-freedom, not two. Two-sided sets fill the interior of the
unit square; the exact set is the origin.

The perfect filter line is not a passive labeling. Chapter~6
(\Cref{ch:set-composition}) will derive the union, intersection, and
complement of two Bernoulli sets and show that
\begin{itemize}
  \item the union of two FNR-zero sets is FNR-zero;
  \item the intersection of two FNR-zero sets is FNR-zero;
  \item the complement swaps the two axes, sending an FNR-zero set
    to an FPR-zero set and vice versa.
\end{itemize}
The FPR-zero side is in general not closed under the same operations:
the union of two FPR-zero sets is FPR-zero only when both are exact
on the same designated set, a fragile condition. The asymmetry
parallels the construction-side asymmetry: FNR-zero is the natural
default in many places, FPR-zero requires more careful design.
\Cref{sec:bridge-to-sets} of chapter~3 already previewed the
union-as-bit-OR identity for FNR-zero Bloom filters as a special case
of the composition theorem; chapter~6 generalizes the calculation and
makes the perfect filter line into an algebraic constraint.

Two further chapters take up the line as their concern. Chapter~14
(\Cref{ch:optimal-bernoulli-sets}) develops the Bernoulli Hash
Function and shows that within the FNR-zero family the bit-per-item
cost is bounded below by an entropy formula, and the BHF achieves the
bound. Chapter~15 (\Cref{ch:perfect-hash-filters}) takes the FPR-zero
side seriously and constructs the perfect hash filter: an approximate
set with $\fprate = 0$ on a designated set $A$ and $\fprate > 0$
outside. Each chapter is one half of the same picture, sitting on one
half of the perfect filter line.

\subsection*{Looking ahead}

We have separated the approximate sets of Part~II into three rate
shapes: FNR-zero (the Bloom filter and its kin), FPR-zero (the perfect
hash filter and its kin), and two-sided (counting filters and
sketches). The shapes are not interchangeable. Each carries different
safety guarantees, different space costs, and different behavior under
algebraic operations on sets. \Cref{sec:ch5-what-comes-next} surveys
the trajectory of the rest of the chapter and of Part~II as a whole:
the composition algebra of chapter~6, the classification measures of
chapter~7, the entropy bound of chapter~8, and the space-optimal
constructions of chapter~14 that close the gap the standard Bloom
filter leaves on the table.
